\begin{ouverture}[algorithmique]{5}{Égalités, programmes et suites}{Partie 1 · Nombres entiers et calculs}
{Nombres cachés · Programmes de calcul · Motifs}
\objectif{Compléter des égalités et trouver des nombres cachés.}
\objectif{Exécuter et produire des programmes de calcul.}
\objectif{Prolonger des suites et calculer une étape lointaine.}
\prerequis{Les quatre opérations et le sens du signe égal.}
\end{ouverture}
\pagestyle{otMain}\essentiel\fontsize{11}{15}\selectfont
\phantomsection\label{cours-5-1}
\essSection{1}{Des égalités et des programmes}
\begin{definition}Le signe $=$ relie deux écritures du même nombre : $178-\square=6\times8$ se résout en calculant d’abord $6\times8=48$, puis $\square=178-48=130$.\end{definition}
\label{fig-cours-5-1-prog-cours}\figureGeo{prog-cours}

\begin{exemple}Exécuter le programme avec 6 : $6+2=8$ ; $8\times4=32$ ; $32-3=29$. Avec 10 : $12\times4=48$ ; $48-3=45$. Pour remonter un programme, on défait chaque étape en sens inverse : si le résultat est 29, on calcule $29+3=32$, puis $32:4=8$, puis $8-2=6$.\end{exemple}
\begin{savoirfaire}[Inventer un programme]Produire un programme, c’est choisir des étapes qui transforment l’entrée en la sortie voulue : passer de 5 à 26 peut s’écrire « multiplier par 5, ajouter 1 ». Tester ensuite son programme sur plusieurs nombres.\end{savoirfaire}
\begin{attention}L’ordre des étapes compte : « ajouter 2 puis multiplier par 4 » transforme 6 en 32 ; « multiplier par 4 puis ajouter 2 » transforme 6 en 26.\end{attention}
\clearpage
\phantomsection\label{cours-5-2}
\essSection{2}{Des suites à prolonger}
\begin{propriete}Une suite suit une règle. Additive : $7 ;\ 15 ;\ 23$ (on ajoute 8). Multiplicative : $7 ;\ 14 ;\ 28$ (on double). Une règle peut combiner : dans $7 ;\ 15 ;\ 31 ;\ 63$, chaque terme est le double du précédent plus 1.\end{propriete}
\begin{savoirfaire}[Atteindre une étape lointaine]Dans une suite de motifs, compter ce qui change à chaque étape. Si l’étape 1 compte 5 carreaux et que chaque étape ajoute 4 carreaux, l’étape 10 compte $5+9\times4=41$ carreaux : neuf ajouts après la première étape.\end{savoirfaire}
\begin{exemple}Étape inverse : quelle étape utilise 33 carreaux ? $33-5=28$, soit $28:4=7$ ajouts : c’est l’étape 8. Le calcul remplace le dessin de toutes les étapes.\end{exemple}
\begin{attention}Toujours vérifier la règle sur plusieurs termes : $2 ;\ 4$ peut continuer en $6$ (ajouter 2) ou en $8$ (doubler). Deux termes ne suffisent pas à fixer une règle.\end{attention}

\clearpage\rubrique{Découverte}
\begin{blocLivre}{ex-5-1}\begin{exercice}[1]{Égalités à trous}
Compléter : $178-\square=6\times8$ ; $\square:4=25$ ; $9\times\square=630$.
\end{exercice}\end{blocLivre}
\begin{blocLivre}{ex-5-2}\begin{exercice}[1]{Exécuter le programme}
Exécuter avec 6, puis avec 10 : choisir un nombre ; ajouter 2 ; multiplier par 4 ; retirer 3.
\label{fig-ex-5-2}\figureGeo{prog-2}
\end{exercice}\end{blocLivre}
\begin{blocLivre}{ex-5-3}\begin{exercice}[1]{Remonter le programme}
Le programme « multiplier par 3 ; ajouter 6 ; multiplier par 2 » affiche 48. Quel nombre a été choisi ?
\end{exercice}\end{blocLivre}
\begin{blocLivre}{ex-5-4}\begin{exercice}[1]{Deux programmes, deux sorties}
Comparer les sorties des programmes « ajouter 2 puis multiplier par 4 » et « multiplier par 4 puis ajouter 2 », pour l’entrée 6, puis pour l’entrée 10.
\end{exercice}\end{blocLivre}
\begin{blocLivre}{ex-5-5}\begin{exercice}[1]{La règle de la suite}
Donner la règle et les deux termes suivants : $7 ;\ 15 ;\ 23 ;\ 31 ;\ \ldots$, puis $5 ;\ 10 ;\ 20 ;\ 40 ;\ \ldots$
\end{exercice}\end{blocLivre}

\rubrique{Entraînement}
\begin{blocLivre}{ex-5-6}\begin{exercice}[2]{Double plus un}
Vérifier que la suite $7 ;\ 15 ;\ 31 ;\ 63$ suit la règle « doubler puis ajouter 1 », puis donner les deux termes suivants.
\end{exercice}\end{blocLivre}
\begin{blocLivre}{ex-5-7}\begin{exercice}[2]{Le nombre caché}
Trouver le nombre caché : $\square\times6-7=41$.
\end{exercice}\end{blocLivre}
\begin{blocLivre}{ex-5-8}\begin{exercice}[2]{Produire un programme}
Écrire un programme de calcul en deux étapes qui transforme 4 en 23, puis un autre, différent, qui transforme aussi 4 en 23.
\end{exercice}\end{blocLivre}
\begin{blocLivre}{ex-5-9}\begin{exercice}[2]{La suite qui zigzague}
Donner les trois termes suivants : $50 ;\ 59 ;\ 56 ;\ 65 ;\ 62 ;\ 71 ;\ \ldots$
\end{exercice}\end{blocLivre}
\begin{blocLivre}{ex-5-10}\begin{exercice}[2]{Des symboles cohérents}
Les mêmes symboles cachent les mêmes nombres : $\triangle+\triangle+\triangle=27$ et $\triangle+\bullet=17$. Trouver les deux nombres.
\end{exercice}\end{blocLivre}

\rubrique{Problèmes}
\begin{blocLivre}{ex-5-11}\begin{exercice}[3]{Le motif qui grandit}
Un motif compte 5 carreaux à l’étape 1 et gagne 4 carreaux à chaque étape.\begin{questions}
\item Combien de carreaux à l’étape 10 ?
\item À quelle étape le motif compte-t-il exactement 33 carreaux ?
\end{questions}
\label{fig-ex-5-11}\figureGeo{motifs-11}
\end{exercice}\end{blocLivre}
\begin{blocLivre}{ex-5-12}\begin{exercice}[3]{Les tee-shirts, suite}
Un tee-shirt coûte 12 euros, la livraison 5 euros quel que soit le nombre acheté. Le prix pour $N$ tee-shirts s’écrit $(N\times12)+5$. Calculer le prix pour 7 tee-shirts, puis trouver $N$ pour un prix de 125 euros.
\end{exercice}\end{blocLivre}
\begin{blocLivre}{ex-5-13}\begin{exercice}[3]{Le tableau de la machine}
Une machine applique le même programme à chaque entrée : $3\to17$ ; $5\to27$ ; $8\to42$. Trouver le programme, puis la sortie pour 20.\begin{center}\renewcommand{\arraystretch}{1.35}\begin{tabular}{c|cccc}\toprule
Entrée & 3 & 5 & 8 & 20\\\midrule
Sortie & 17 & 27 & 42 & ?\\
\bottomrule\end{tabular}\end{center}
\end{exercice}\end{blocLivre}
\begin{blocLivre}{ex-5-14}\begin{exercice}[3]{Vrai ou faux ?}
Nora affirme : « La suite $2 ;\ 4 ;\ \ldots$ continue forcément par 6. » Donner deux suites différentes qui commencent par 2 puis 4, avec leurs règles.
\end{exercice}\end{blocLivre}
\begin{blocLivre}{ex-5-15}\begin{exercice}[3]{Le défi du professeur}
Un professeur écrit : « $\square\times\square=36$ et $\square$ est pair. » Trouver toutes les valeurs possibles du nombre caché, sachant qu’un même symbole cache le même nombre.
\end{exercice}\end{blocLivre}

\clearpage
